\subsection{Automatic pivotality and tensor-quotient rigidity}
\label{subsec:frontier-tensor-rigidity}
The pivotal compatibility assumed in Theorem~\ref{thm:r1660-pivotal-quotient}
is automatic after passage to the positive Galois form and its dominant
image. We give the general argument before applying it to both residual
rings.

\begin{lemma}[Pseudo-unitary images and canonical pivotal structures]
\label{lem:pu-tensor-functors}
Let $F\colon\mathcal C\to\mathcal D$ be a tensor functor between
complex fusion categories.
\begin{enumerate}
\item If $F$ is dominant and $\mathcal C$ is pseudo-unitary, then
$\mathcal D$ is pseudo-unitary.
\item If $\mathcal C$ and $\mathcal D$ are pseudo-unitary, then $F$
preserves their canonical spherical structures.
\end{enumerate}
\end{lemma}

\begin{proof}
For (i), regard $\mathcal D$ as a left $\mathcal C$-module category via
$X\cdot M=F(X)\otimes M$. Dominance implies that this module category
is indecomposable: its component containing $\mathbf1$ contains every
simple summand of every $F(X)$, hence every simple object of $\mathcal D$.
Its dual $\mathcal E=\operatorname{Fun}_{\mathcal C}
(\mathcal D,\mathcal D)$ is therefore a fusion category.
It has the same categorical and Frobenius--Perron dimensions as
$\mathcal C$, by \cite[Corollaries~9.3.3 and~7.16.7 and Proposition~9.3.9]{EGNO}. Thus $\mathcal E$ is pseudo-unitary.

Right tensor multiplication defines a fully faithful tensor functor
$\mathcal D^{\mathrm{rev}}\to\mathcal E$ (with composition as the
tensor product in $\mathcal E$). Indeed, a $\mathcal C$-module natural
transformation between $-\otimes Y$ and $-\otimes Z$ has component
$\mathrm{id}_{F(X)}\otimes f$ at $F(X)$, where its component at
$\mathbf1$ is $f\colon Y\to Z$. For a direct summand $M$ of $F(X)$,
naturality with the splitting maps forces its component to be
$\mathrm{id}_M\otimes f$. Dominance proves the assertion for every
$M$, and conversely every $f$ gives such a transformation.
Restricting the positive spherical structure of $\mathcal E$ to this
fusion subcategory gives a positive spherical structure on
$\mathcal D^{\mathrm{rev}}$, and hence on $\mathcal D$. Its dimensions
equal Frobenius--Perron dimensions, so $\mathcal D$ is pseudo-unitary.

We first record a finite-spectrum observation for (ii). Every monoidal
natural automorphism $u$ of a tensor functor between fusion categories
has finite order. To see this, choose representatives $X_i$ of the
finitely many source simples. The union of the spectra of $u_{X_i}$,
and the maximal size of their Jordan blocks, are finite. Naturality
with direct-sum decompositions gives the same spectrum set and Jordan
bound for $u_X$ for every source object $X$. If $\lambda$ is an
eigenvalue of $u_X$, a nonzero eigensubobject has nonzero tensor powers,
so monoidality shows that $\lambda^n$ is an eigenvalue of $u_{X^{\otimes n}}$
for every $n\geq1$. Thus $\lambda$ is a root of unity.

There cannot be a Jordan block of size $r>1$. Such a block occurs on
a multiplicity space of a simple target object $Y$, and its $n$th
tensor power occurs on a multiplicity space tensored with the nonzero
object $Y^{\otimes n}$. The $n$th tensor power of a size-$r$ Jordan
block with nonzero eigenvalue has a block of size $n(r-1)+1$ in
characteristic zero. Explicitly, after removing its eigenvalue, write
the block as $\exp(N)$, with $N^{r-1}\ne0$ and $N^r=0$.
Its tensor power is the exponential of the sum of $n$ commuting
copies of $N$; the power $n(r-1)$ of that sum is nonzero, by the
multinomial formula, and its next power vanishes. This contradicts
the uniform Jordan bound. Therefore all the $u_{X_i}$ are diagonalizable
with root-of-unity eigenvalues. A common multiple of their eigenvalue
orders is the order of $u$.

Now write $j_{\mathcal C}$ and $j_{\mathcal D}$ for the canonical
spherical structures, whose existence and uniqueness are
\cite[Proposition~9.5.1]{EGNO}. The canonical comparison between $F$ and
double duals identifies $F(j_{\mathcal C,X})$ with a monoidal natural
isomorphism $F(X)\to F(X)^{**}$. Comparing this with
$j_{\mathcal D,F(X)}$ gives a monoidal natural automorphism
$u$ of $F$. Preservation of evaluation and coevaluation by $F$ gives
\[
 \operatorname{tr}_{\mathcal D}(u_X)
 =\dim_{\mathcal C}(X)
 =\operatorname{FPdim}(X)
 =\operatorname{FPdim}(F(X))
 =\dim_{\mathcal D}(F(X)).
\]
Here preservation of Frobenius--Perron dimension is
\cite[Proposition~4.5.7]{EGNO}. By the observation just proved, $u_X$ is
diagonalizable with eigenvalues of modulus one. Its categorical trace
is the sum of these eigenvalues weighted by the positive dimensions
of the corresponding eigensubobjects. The displayed equality forces
every eigenvalue to be $1$. Hence $u_X=\mathrm{id}_{F(X)}$ for every
$X$, which is exactly pivotal compatibility.
\end{proof}

\begin{theorem}[Tensor-quotient rigidity of the residual rings]
\label{thm:frontier-all-tensor-functors}
Let $\mathcal C$ be a complex fusion category with Grothendieck ring
$\RR{4}{155}$ or $\RR{8}{1660}$. Every tensor functor from
$\mathcal C$ to a complex fusion category is fully faithful.
Equivalently, every dominant tensor functor from $\mathcal C$ is a
tensor equivalence.
\end{theorem}

\begin{proof}
First suppose $F\colon\mathcal C\to\mathcal D$ is dominant.
Fusion categories and tensor functors between them are defined over
algebraic number fields \cite[Corollary~9.1.8]{EGNO}; choose a common
such field. Extend the Galois conjugation of
Proposition~\ref{prop:nonselfdual-neargroup-spherical} for $\RR{4}{155}$,
or the sphericalization argument in Section~\ref{subsec:frontier-1660}
for $\RR{8}{1660}$, to this field. It transports
$F$ to a dominant tensor functor $F^\tau\colon\mathcal C^\tau\to
\mathcal D^\tau$, with $\mathcal C^\tau$ pseudo-unitary and with the
same based ring. Lemma~\ref{lem:pu-tensor-functors} makes
$\mathcal D^\tau$ pseudo-unitary and $F^\tau$ pivotal for the canonical
structures. An equivalence after conjugation is an equivalence before
conjugation. Thus we may work in the positive spherical form.

For $\RR{8}{1660}$, Theorem~\ref{thm:r1660-pivotal-quotient} now makes
$F$ an equivalence.

It remains to consider $\RR{4}{155}$. Write
$U:\mathcal Z(\mathcal C)\to\mathcal C$ for the forgetful functor and
use the established decomposition
\[
 \mathcal Z(\mathcal C)\simeq\mathcal S\boxtimes\mathcal B,
 \qquad \theta_s=i,\qquad U(s\boxtimes\mathbf1)=g,
\]
where $s$ is the nontrivial semion. The basis of $K_0(\mathcal C)$ is
$\mathbf1,h,g,hg,X,Y$, with $gX=Y$ and $hX=X$.
For clarity denote the simple of $\mathcal B$ with numerical label $j$
by $B_j$; thus $B_1=\mathbf1$.
The complete list of twist-one simples of the center, with their
forgetful images, is
\begin{equation}\label{eq:r155-twist-one-restrictions}
\begin{array}{c|l}
\text{central simple}&\text{forgetful image}\\ \hline
 \mathbf1\boxtimes B_1&\mathbf1\\
 \mathbf1\boxtimes B_2&Y\\
 \mathbf1\boxtimes B_5&X\\
 \mathbf1\boxtimes B_{23}&\mathbf1+2X+2Y\\
 \mathbf1\boxtimes B_{25}&\mathbf1+h+2X+2Y\\
 s\boxtimes B_{18},\ s\boxtimes B_{19}&\mathbf1+X+Y\\
 s\boxtimes B_{24}&\mathbf1+h+2X+2Y.
\end{array}
\end{equation}
Indeed the twist-one labels in $\mathcal B$ are $1,2,5,23,25$, while
the twist-$(-i)$ labels are $18,19,24$. Before tensoring by $s$, the
last three restrict respectively to $g+X+Y$, $g+X+Y$, and
$g+hg+2X+2Y$.

Let $R$ be the right adjoint of $F$.
Its induced central algebra $E$ has underlying algebra $R(\mathbf1)$,
and
\[
 \dim\operatorname{Hom}_{\mathcal C}(\mathbf1,U(E))
 =\dim\operatorname{Hom}_{\mathcal D}(F(\mathbf1),\mathbf1)=1.
\]
The induced central-algebra construction of Brugui\`eres--Natale
\cite[Section~5]{BN} identifies $\mathcal D$ with the category of modules over $U(E)$.
In characteristic zero this induced central algebra is separable:
the underlying module category is semisimple, and the bimodule
splitting of multiplication can be averaged to a morphism in the
center. Thus $E$ is connected \'etale in $\mathcal Z(\mathcal C)$.
Positivity of the spherical structure gives $\theta_E=\mathrm{id}_E$;
see Schopieray \cite[Lemma~2.2.4]{SchLevelBounds}.
The displayed connectedness in $\mathcal C$, together with
\eqref{eq:r155-twist-one-restrictions}, therefore leaves only
\[
 E\simeq\mathbf1\oplus n(\mathbf1\boxtimes B_2)
                  \oplus n(\mathbf1\boxtimes B_5),\qquad n\geq0.
\]
Here the two multiplicities agree because $B_2^*=B_5$ and an \'etale
algebra has a nondegenerate Frobenius pairing.

Put $Z=\mathbf1\boxtimes B_2$. The reconstructed fusion rules give
\[
 B_2^2=B_9+B_{10}+B_{16}+B_{19},\qquad
 B_2B_5=B_1+B_8+B_{11}+B_{22}.
\]
Thus multiplication on any two copies of $Z$ in $E$ vanishes, and
multiplication on a copy of $Z$ and a copy of $Z^*$ has image in the
unit summand. If $n>0$, nondegeneracy of the Frobenius pairing supplies
copies $Z\hookrightarrow E$ and $Z^*\hookrightarrow E$ for which this
latter multiplication is nonzero. Associativity on
$Z\otimes Z\otimes Z^*$ is then impossible: multiplying the first
two factors gives zero, whereas multiplying the last two and then the
first gives a nonzero multiple of
$\mathrm{id}_Z\otimes\mathrm{ev}_Z$. Consequently $n=0$ and
$R(\mathbf1)\simeq\mathbf1$. The projection formula gives
$RF(V)\simeq V\otimes R(\mathbf1)\simeq V$; hence $F$ is fully
faithful and, by dominance, an equivalence.


Finally, an arbitrary tensor functor factors through the full fusion
subcategory generated by its image. Its first factor is dominant and
hence, by what we have proved, an equivalence. Its second factor is a
fully faithful inclusion. This proves the statement for every tensor
functor.
\end{proof}

This theorem gives a restriction on hypothetical categorifications;
it neither constructs nor excludes either one. The two general
categorification questions remain open.
